\documentclass[12pt]{article}
\usepackage[utf8]{inputenc}
\usepackage[T1]{fontenc}
\usepackage{hyperref}
\title{Woiwurrung Language Resources for K-12 Curriculum Development}
\author{Compiled by Claude Sonnet 5.5 (claude-sonnet-5-5) for Andr\'as Kornai}
\date{}
\begin{document}
\maketitle

\section{Introduction}
Woiwurrung (Woi-wurrung; ISO 639-3 wyi; Glottocode woiw1237) is a Pama-Nyungan language of the Melbourne and Yarra River area of Victoria, Australia. It is the language of the Wurundjeri people. It is dormant, but community revival work is under way. The searches did not turn up a dedicated modern Woiwurrung dictionary or school curriculum.

\section{Primary Educational Resources}
\begin{itemize}
\item Woiwurrung early years kit (45 language cards and an audio CD, with artwork and recordings by Mandy Nicholson), produced by Yarra Ranges Council, HICSA and Wurun Child and Family Place; reported in a local news article: \url{https://upperyarra.mailcommunity.com.au/mail/2015-10-09/indigenous-words-launch/}
\item Woi-wurrung language session for children, run by Wurundjeri educators (greetings, numbers, animals, body parts), Moonee Valley library listing: \url{https://libraryevents.mvcc.vic.gov.au/event?id=226295}
\end{itemize}

\section{Linguistic and Reference Materials}
\begin{itemize}
\item AUSTLANG entry S36 Woiwurrung, with alternative names and spellings: \url{https://collection.aiatsis.gov.au/austlang/language/s36}
\item Victorian Aboriginal Corporation for Languages (VACL) items catalogued in Victorian Collections, including the \emph{Wurrung Wurrung Victorian Aboriginal languages resource kit} (2008) and the \emph{Dictionary of Aboriginal placenames of Melbourne and Central Victoria} (2002): \url{https://victoriancollections.net.au/collections?q=woiwuru}
\item Glottolog reference record for Clark and Heydon (2002): \url{https://glottolog.org/resource/reference/id/5428}
\end{itemize}

\section{Religious and Community Resources}
\begin{itemize}
\item Report on the exhibition \emph{dhumba-njan dhumba-njarr}, curated by Brooke Wandin to revive the Wurundjeri language: \url{https://upperyarra.mailcommunity.com.au/entertainment/2019-05-14/exhibition-hopes-to-revive-wurundjeri-language/}
\item Woiwurrung language awareness session with Djirri Djirri (Wurundjeri women's dance group), Djirra: \url{https://djirra.org.au/?p=6911}
\end{itemize}
I saw no scripture resources in the search results.

\section{Digital and Audio Resources}
The early years kit above includes an audio CD. I found no online dictionary or audio archive in the results.

\section{Curriculum Development Context}
The language is dormant. Revival is led by the Wurundjeri Woi-wurrung Cultural Heritage Aboriginal Corporation and Wurundjeri educators, who run children's and adult sessions. A school curriculum would depend on historical wordlists, the VACL materials and Wurundjeri direction. Reconstructed grammar and vocabulary would need community approval. I did not confirm the current programmes of the Wurundjeri corporation, so contact the corporation and VACL directly.

\end{document}